\clearpage\fancyhead[L]{\small GAUSS TO VON LINDENAU.}\SourceLeaf{321}

\begin{center}\textbf{GAUSS TO VON LINDENAU.}\par\vspace{4pt}\hrule\vspace{4pt}\fontsize{10}{12}\selectfont Monatliche Correspondenz. Vol. XXI. pp. 276\,..\,280. March 1810.\par\vspace{4pt}\hrule\vspace{4pt}\end{center}
\begin{flushright}Göttingen, 23 Feb. 1810.\end{flushright}

\par\indent ---\enspace---\enspace--- I have just finished some calculations on the last two oppositions of Pallas. I had the results printed in no. 32 of the Göttinger gelehrte Zeitungen [p. 317 of this volume], and I am sending them to you here. ---\enspace---

I have already derived some further results about the elements from these calculations, and I am now working on a few more calculations on the subject. I plan to make this into a paper for our Society, along with an account of several techniques I have used for many years in applying the method (\emph{des moindres carrés}) developed in Section 3 of Book 2 of my \emph{Theoria}. I will send you an excerpt as soon as the work is finished.

I am including our local observations of \emph{Vesta}, all made by Professor \textsc{Harding} with the mural quadrant.

\Needspace{12\baselineskip}\begin{center}\emph{Observations of Vesta.}\end{center}
\begin{center}\fontsize{10}{12}\selectfont\setlength{\tabcolsep}{5pt}\renewcommand{\arraystretch}{1.15}\begin{tabular}{lrrr}
1810 & \shortstack{Mean time\\in Göttingen} & \shortstack{Apparent R.A.\\of Vesta} & \shortstack{Apparent northern\\declination}\\
January 13 & $11^{\text{u}}\;1^{\text{m}}\;27^{\text{s}}7$ & $98^\circ\;5'\;44''5$ & $23^\circ\;19'\;17''6$\\
15 & $10\;51\;30.1$ & $97\;34\;57.0$ & $23\;26\;14.3$\\
28 & $9\;49\;13.7$ & $94\;46\;15.2$ & $24\;7\;33.7$\\
February 2 & $9\;26\;32.7$ & $94\;0\;51.3$ & $24\;20\;51.7$\\
3 & $9\;22\;7.4$ & $93\;53\;19.2$ & $24\;23\;37.8$\\
6 & $9\;8\;56.1$ & $93\;32\;33.2$ & $24\;31\;1.9$\\
\end{tabular}\end{center}

\Needspace{21\baselineskip}\par\indent Here are also our observations of the occultation of Jupiter on 8 February. The emersions are quite uncertain because the Moon was very low.

\Needspace{17\baselineskip}\begin{center}\emph{Immersions.}\end{center}
\begin{center}\fontsize{10}{12}\selectfont\setlength{\tabcolsep}{5pt}\renewcommand{\arraystretch}{1.0}\begin{tabular}{lrl}
1. satellite & $8^{\text{u}}\;58^{\text{m}}\;2^{\text{s}}9$ & \textsc{Harding}\\
 & $3.3$ & \textsc{Gauss}\\
2. satellite & $8\;58\;26.2$ & G.\\
 & $28.7$ & H.\\
$\JupiterSymbol$ I. R. & $9\;0\;16.4$ & H.\\
 & $17.4$ & G.\\
$\JupiterSymbol$ II. R. & $9\;2\;26.0$ & G.\\
 & $26.7$ & H.\\
3. satellite & $9\;7\;55.5$ & G.\\
 & $56.4$ & H.\\
4. satellite & $9\;14\;35.3$ & H.\\
 & $37.0$ & G.\\
\end{tabular}\end{center}
\Needspace{9\baselineskip}\begin{center}\emph{Emersions.}\end{center}
\begin{center}\fontsize{10}{12}\selectfont\setlength{\tabcolsep}{5pt}\renewcommand{\arraystretch}{1.0}\begin{tabular}{lrl}
1. satellite & $9^{\text{u}}\;30^{\text{m}}\;59^{\text{s}}$ & \textsc{Gauss} already emerged\\
$\JupiterSymbol$ I. R. & $9\;33\;37.9$ & H.\\
II. R. & $9\;35\;19.6$ & G.\\
 & $36.1$ & H.\\
3. satellite & $9\;43\;39.3$ & H.\\
\end{tabular}\end{center}

\par\indent I am much obliged to Mr. \textsc{Oriani} for noting the misprints in my Theoria. He is quite right that on p. 129 I forgot to add that $B, B', B'' = 0$ must be assumed
\SourceLeaf{322} if the condition equation under which equation [7] is unusable is to take the form given there. It is clear, moreover, that even when $B, B', B''=0$ does not hold, equation [7] can still be unusable if the twelve-term expression that \textsc{Oriani} developed happens to become $=0$ or very small. It had already occurred to me that \textsc{Euler} had found the theorem from which the fine theorem I attributed to \textsc{La Place} can be derived very easily. By then, however, the passage on p. 212 had already been printed. I chose to leave it out of the errata, because at least \textsc{La Place} was the first to state that theorem in the form used there. I had already noted most of the misprints reported by \textsc{Oriani}. Here are three others he missed:

\begin{center}\fontsize{10}{12}\selectfont\setlength{\tabcolsep}{5pt}\renewcommand{\arraystretch}{1.15}\begin{tabular}{lrlrlllll}
P. & 1 & Line & 4 & from bottom & for & inversa & read & composita\\
--- & 65 & --- & 2 & --- & --- & $\cos\theta$. & read & $\cos\zeta$\\
--- & 195 & --- & 16 &  & --- & $\lambda'''$ & read & $=\lambda'''$.\\
\end{tabular}\end{center}

\par\indent I cannot find the misprint $380^\circ$ for $180^\circ$ on p. 4 in my copy. In place of the two misprints on p. 83, one could give a single correction, namely $\zeta+45$ in place of $\zeta-45^\circ$.
